\subsection{Secondary sensitivity results}
\label{or:aor-secondary-comparative-statics}

This section retains sensitivity calculations that are useful for interpreting
the economic model but are not needed for the safe-compression reduction.
Each calculation holds along the named parameter path or conditional on the
displayed initial belief and fixed action. Unnamed primitives are fixed.
Finiteness is used for termwise differentiation and finite action
maximization. No closure-detectability or stationary-contraction argument is
used.

\subsection{Bridge-margin sensitivity}

Write
\begin{equation}
 A_{\mathrm{br}}:=\beta^d\rho^d\pi C,
 \qquad M_{\mathrm{br}}:=A_{\mathrm{br}}-\kappa,
 \qquad m_{\mathrm{br}}:=\frac{M_{\mathrm{br}}}{A_{\mathrm{br}}}.
 \label{eq:or-bridge-margin}
\end{equation}
On the declared real extension with fixed integer \(d\), positive gross
coordinates, and \(M_{\mathrm{br}}>0\), realized bridge value equals the
margin. Its named-coordinate elasticities are
\begin{equation}
 (\varepsilon_\beta^M,\varepsilon_\rho^M,
   \varepsilon_\pi^M,\varepsilon_C^M,\varepsilon_\kappa^M)
 =\left(\frac d{m_{\mathrm{br}}},\frac d{m_{\mathrm{br}}},
   \frac1{m_{\mathrm{br}}},\frac1{m_{\mathrm{br}}},
   -\frac{1-m_{\mathrm{br}}}{m_{\mathrm{br}}}\right).
 \label{eq:or-bridge-elasticities}
\end{equation}
Threshold fragility is governed by the normalized positive margin
\(m_{\mathrm{br}}\). The magnitudes diverge as
\(m_{\mathrm{br}}\downarrow0\), but not merely because
\(M_{\mathrm{br}}\downarrow0\): if cost is zero and gross scale vanishes,
then \(m_{\mathrm{br}}=1\). At \(M_{\mathrm{br}}=0\), bridge-value
elasticity is undefined. At and below that threshold, the signed margin,
action region, or exact finite changes in realized value replace a percentage
elasticity.

For the derivation, condition on the canonical project and hold every unnamed
coordinate fixed. Differentiating on the positive real extension gives
\[
 \beta\,\partial_\beta M_{\mathrm{br}}=dA_{\mathrm{br}},
 \quad
 \rho\,\partial_\rho M_{\mathrm{br}}=dA_{\mathrm{br}},
 \quad
 \pi\,\partial_\pi M_{\mathrm{br}}
 =C\,\partial_C M_{\mathrm{br}}=A_{\mathrm{br}},
 \quad
 \kappa\,\partial_\kappa M_{\mathrm{br}}=-\kappa.
\]
On \(M_{\mathrm{br}}>0\), realized value
\([M_{\mathrm{br}}]_+\) equals the signed margin. Dividing by
\(M_{\mathrm{br}}\) and using
\(m_{\mathrm{br}}=M_{\mathrm{br}}/A_{\mathrm{br}}\) gives
\eqref{eq:or-bridge-elasticities}. At the threshold, the positive-part value
has no ordinary percentage elasticity. The costless path
\(\kappa=0\), \(A_{\mathrm{br}}\downarrow0\) keeps
\(m_{\mathrm{br}}=1\), so a vanishing net level alone does not imply the
threshold-fragility limit.

\subsection{Operational and generative elasticity contributions}

Let one named positive scalar \(x\) vary the same primitive in total insertion
value \(I(x)\), operational value \(O(x)\), and generative value \(G(x)\),
holding every unlisted primitive fixed and remaining on a differentiable
branch. Differentiating \(I=O+G\) gives
\[
 xI'(x)=xO'(x)+xG'(x).
\]

\begin{proposition}[Operational--generative elasticity contributions]
\label{prop:or-channel-elasticity}
At \(x_0>0\), suppose \(I(x)=O(x)+G(x)\) throughout a neighborhood of
\(x_0\), and the three channel paths are differentiable there along the same
named perturbation. If \(I(x_0)>0\), define
\[
 C_x^{\mathrm{op}}:=\frac{x_0O'(x_0)}{I(x_0)},
 \qquad
 C_x^{\mathrm{gen}}:=\frac{x_0G'(x_0)}{I(x_0)}.
\]
Then
\[
 \varepsilon_x^I=C_x^{\mathrm{op}}+C_x^{\mathrm{gen}}.
\]
Only when both channel levels are positive may this be written as
\[
 \varepsilon_x^I
 =\frac{O}{I}\varepsilon_x^{\mathrm{op}}
  +\frac{G}{I}\varepsilon_x^{\mathrm{gen}}.
\]
\end{proposition}

The contribution form remains meaningful when one channel level is zero; the
corresponding ordinary component log elasticity does not. If total value is
zero, no total elasticity or total-normalized contribution is reported. The
scaled derivative identity above, or the exact finite-change identity
\(\Delta I=\Delta O+\Delta G\), remains well defined.

Along one fixed differentiable scalar path, differentiating the level identity
and applying common scaling gives
\[
 xI'(x)=xO'(x)+xG'(x).
\]
If \(I(x)>0\), division by the total level gives the two signed contributions
in \cref{prop:or-channel-elasticity}. Rewriting a contribution as
\[
 \frac{xO'}I=\frac OI\frac{xO'}O
\]
requires \(O>0\), and the analogous rewrite requires \(G>0\). If a component
vanishes, its scaled derivative remains available but its log elasticity does
not. If \(I=0\), neither total elasticity nor a total-normalized contribution
is defined; the level derivative or exact finite-change identity is used.

\begin{table}[ht]
\centering
\scriptsize
\setlength{\tabcolsep}{2.5pt}
% Generated by julia/scripts/run_unified_resource_benchmark.jl
\begin{tabular}{llrrrrrr}
\toprule
Belief & Library & $V$ & $V^{\mathrm{op}}$ & $V^{\mathrm{gen}}$ & $D_\beta$ & $E^{\mathrm{op}}_\beta$ & $E^{\mathrm{gen}}_\beta$ \\
\midrule
low & I & $\frac{115}{288}$ & $0$ & $\frac{115}{288}$ & $\frac{10777}{2070}$ & $--$ & $\frac{10777}{2070}$ \\
low & I+A & $1$ & $0$ & $1$ & $\frac{56}{15}$ & $--$ & $\frac{56}{15}$ \\
low & I+D & $\frac{14}{3}$ & $\frac{14}{3}$ & $0$ & $\frac{25}{21}$ & $\frac{25}{21}$ & $--$ \\
high & I & $\frac{23}{288}$ & $0$ & $\frac{23}{288}$ & $\frac{14089}{2070}$ & $--$ & $\frac{14089}{2070}$ \\
high & I+A & $\frac{7}{6}$ & $0$ & $\frac{7}{6}$ & $\frac{53}{15}$ & $--$ & $\frac{53}{15}$ \\
high & I+D & $\frac{22}{3}$ & $\frac{22}{3}$ & $0$ & $\frac{29}{33}$ & $\frac{29}{33}$ & $--$ \\
\bottomrule
\end{tabular}

\caption{Exact productive value, channel levels, and fixed-policy discount
elasticity in the canonical benchmark. A dash marks an undefined ordinary
component elasticity at a zero component level; it is never recorded as
zero. Scaled contributions are evaluated directly.}
\label{tab:or-canonical-channel-elasticities}
\end{table}

\subsection{Fixed-exposure innovation duration}

For a fixed nonnegative, nonzero exposure sequence
\(z_0,\ldots,z_{H-1}\) and effective discount \(\alpha>0\), let
\begin{equation}
 \Psi_H(\alpha):=\sum_{t=0}^{H-1}\alpha^tz_t,
 \quad \omega_t:=\frac{\alpha^tz_t}{\Psi_H(\alpha)},
 \quad D_\Psi:=\sum_t t\omega_t,
 \quad C_\Psi:=\sum_t\omega_t(t-D_\Psi)^2.
 \label{eq:or-innovation-potential-duration}
\end{equation}
Because the sequence is nonzero, \(\Psi_H(\alpha)>0\). The exact finite
scaled-derivative identities are
\begin{equation}
 \alpha\Psi_H'(\alpha)=D_\Psi\Psi_H(\alpha),
 \qquad
 \alpha D_\Psi'(\alpha)=C_\Psi\geq0.
 \label{eq:or-innovation-duration}
\end{equation}
Equivalently, \(D_\Psi\) is the derivative of \(\log\Psi_H\) with respect to
\(\log\alpha\), and \(C_\Psi\) is the corresponding derivative of
\(D_\Psi\). These are interpretations of the displayed scaled identities.
Innovation duration is the contribution-weighted date at which fixed
uncovered value arrives, and innovation convexity is its timing variance. A
higher effective discount shifts weight toward later exposures and weakly
raises duration. This is not primitive project delay, time to admission, or a
convexity theorem for an optimized Bellman value whose exposure sequence or
selected action changes. If exposure is zero, duration and log elasticity are
undefined; level contributions or finite differences are used instead.

For the derivation, introduce
\[
 N_1(\alpha):=\sum_t t\alpha^tz_t,
 \qquad
 N_2(\alpha):=\sum_t t^2\alpha^tz_t.
\]
Termwise finite differentiation gives
\[
 \alpha\Psi_H'(\alpha)=N_1(\alpha),
 \qquad
 \alpha N_1'(\alpha)=N_2(\alpha).
\]
Since \(\Psi_H>0\), \(D_\Psi=N_1/\Psi_H\), and the quotient rule yields
\[
 \alpha D_\Psi'
 =\frac{N_2\Psi_H-N_1^2}{\Psi_H^2}
 =\sum_t\omega_t(t-D_\Psi)^2=C_\Psi\geq0.
\]
The nonnegative curvature is curvature of \(\log\Psi_H\) in
\(\log\alpha\) for a fixed exposure sequence. Changing project duration,
belief dynamics, completion laws, or the selected action changes the sequence
and lies outside this derivative.

\subsection{Finite switch and one-sided-change conventions}

Within a strictly dominant price branch, the penalized envelope inherits the
branch slope \(-W(L')\). A pairwise crossing is only a candidate switch; it
becomes a resource-price breakpoint only when both branches are globally
active. At an unequal-burden tie the optimizer is set-valued, so one-sided
slopes or cross-breakpoint finite changes replace a derivative selected by an
unstated tie breaker.

For a finite library family, record the best-versus-second-best objective
margin; for a Bellman node, record the corresponding action-value margin. A
positive margin measures distance from a tie in value space. A zero margin
flags a possible switch, but distinct libraries or actions may generate
identical branches, so it is not sufficient for a change of optimizer. At an
actual switch with positive value, one-sided or arc elasticities may be
reported; at zero value, only level slopes and finite differences are defined.
